\documentclass[a4paper, serif, 10pt]{moderncv}

%% moderncv
% style and color
\moderncvstyle{banking}
\moderncvcolor{black}

%% page layout/margins
\usepackage[top=2.5cm, bottom=2.5cm, left=1.5cm, right=1.5cm]{geometry}

\nopagenumbers{}

%% Babel config for Indonesian (Bahasa Indonesia) language
% ref:
% - https://latex3.github.io/babel/guides/locale-indonesian.html
\usepackage[indonesian]{babel}

%% fontspec
\usepackage{fontspec}

\IfFontExistsTF{Linux Libertine}{
  \setmainfont[Ligatures={TeX, Common}, Numbers=OldStyle]{Linux Libertine}
}{}
\IfFontExistsTF{Linux Libertine O}{
  \setmainfont[Ligatures={TeX, Common}, Numbers=OldStyle]{Linux Libertine O}
}{}

%% CTeX
% CJK support, used to show CJK characters in the resume
%
% - fontset=none: disable builtin fontset but instead set the CJK font manually
% - heading=false: disable ctex heading
% - punct=kaiming: use kaiming punctuations styles for CJK
% - scheme=plain: use plain scheme, do not override \normalsize font size
% - space=auto: space settings for CJK characters
%
% ref:
% - http://ctan.mirrorcatalogs.com/language/chinese/ctex/ctex.pdf
\usepackage[UTF8, heading=false, punct=kaiming, scheme=plain, space=auto]{ctex}

\IfFontExistsTF{Noto Serif CJK SC}{
  \setCJKmainfont{Noto Serif CJK SC}
}{}
\IfFontExistsTF{Noto Sans CJK SC}{
  \setCJKsansfont{Noto Sans CJK SC}
}{}

%% line spacing
\usepackage{setspace}
\setstretch{1.125}

%% URL styling - use normal text instead of monospace
\urlstyle{same}

% Auto-underline all links
% Defer until \begin{document} so hyperref is already loaded
\AtBeginDocument{%
  \let\oldhref\href
  \renewcommand{\href}[2]{\oldhref{#1}{\underline{#2}}}%
}

%% Patch \httplink and \httpslink to handle full URLs
% The original moderncv macros always prepend http:// or https:// to URLs.
% This causes issues when the URL already has a protocol (e.g., https://example.com)
% resulting in malformed URLs like https://https://example.com.
% This patch checks if the URL already contains "://" and uses it directly if so.
% Note: We use \str_set:Nx (x-expansion) to expand macros like \@homepage first.
\ExplSyntaxOn
\str_new:N \g_yamlresume_url_str

\RenewDocumentCommand{\httpslink}{O{}m}{
  \str_set:Nx \g_yamlresume_url_str {#2}
  \str_if_in:NnTF \g_yamlresume_url_str {://}
    { \url{#2} }
    { \url{https://#2} }
}

\RenewDocumentCommand{\httplink}{O{}m}{
  \str_set:Nx \g_yamlresume_url_str {#2}
  \str_if_in:NnTF \g_yamlresume_url_str {://}
    { \url{#2} }
    { \url{http://#2} }
}
\ExplSyntaxOff

\name{Raka Mahendra}{}
\title{Data Scientist | Machine Learning \& Data Analytics}
\phone[mobile]{+62 812-3456-7890}
\email{raka.mahendra@example.com}
\homepage{https://rakamahendra.dev}

\address{Jl. Senopati No. 12, Kebayoran Baru, Jakarta Selatan, DKI Jakarta, Indonesia, 12190}{}{}

\extrainfo{{\small \faLinkedin}\ \href{https://www.linkedin.com/in/rakamahendra}{@rakamahendra} {} {} {} • {} {} {}
{\small \faGithub}\ \href{https://github.com/rakamahendra}{@rakamahendra} {} {} {} • {} {} {}
{\small \faMedium}\ \href{https://medium.com/@rakamahendra}{@rakamahendra}}

\begin{document}

\maketitle

\section{Informasi Dasar}

\cvline{}{Data Scientist dengan pengalaman lebih dari 5 tahun dalam merancang dan
mengimplementasikan solusi machine learning untuk perusahaan e-commerce dan
fintech. Berpengalaman dalam membangun model prediktif, menyusun kerangka
eksperimen A/B, serta menerjemahkan kebutuhan bisnis menjadi analisis data
yang actionable. Mahir menggunakan Python, SQL, dan platform cloud untuk
mengembangkan pipeline data end-to-end. Memiliki kemampuan komunikasi
lintas tim yang baik dan terbiasa bekerja dalam lingkungan Agile.}
\section{Pendidikan}

\cventry{Agu 2013–Jun 2017}
        {Sarjana, Matematika dan Statistika, Nilai: 3.78}
        {Institut Teknologi Bandung}
        {\url{https://www.itb.ac.id}}
        {}
        {\begin{itemize}
\item Mempelajari fondasi statistika, matematika, dan pemrograman untuk analisis data
\item Menyelesaikan tugas akhir tentang pemodelan prediktif risiko kredit menggunakan regresi logistik
\item Aktif sebagai asisten praktikum Statistika selama dua semester
\end{itemize}
\textbf{Mata Kuliah}: Statistika Matematika, Analisis Data Eksploratif, Pemodelan Statistika, Machine Learning, Riset Operasi, Basis Data, Aljabar Linear Terapan, Pengantar Ilmu Komputer}
\section{Pengalaman Kerja}

\cventry{Mar 2022–Sekarang}
        {Data Scientist}
        {PT Digital Nusantara}
        {\url{https://www.digitalnusantara.co.id}}
        {}
        {\begin{itemize}
\item Memimpin pengembangan model machine learning untuk sistem rekomendasi produk, meningkatkan konversi sebesar 18\%
\item Merancang dan mengevaluasi eksperimen A/B untuk menguji strategi personalisasi dan promosi
\item Mengembangkan pipeline data otomatis menggunakan Apache Airflow dan dbt untuk mendukung kebutuhan analitik real-time
\item Berkolaborasi dengan tim produk dan engineering dalam mengidentifikasi peluang bisnis berbasis data
\item Menerapkan model monitoring untuk memantau kualitas prediksi dan mendeteksi data drift
\end{itemize}
\textbf{Kata Kunci}: Machine Learning, Python, SQL, A/B Testing, Apache Airflow, Model Monitoring, BigQuery}

\cventry{Jul 2019–Feb 2022}
        {Data Analyst}
        {PT Finansial Cerdas}
        {\url{https://www.finansialcerdas.id}}
        {}
        {\begin{itemize}
\item Menyusun dashboard analitik dan laporan berkala untuk memantau performa bisnis
\item Melakukan analisis kohort dan segmentasi pelanggan untuk mendukung strategi retensi
\item Menggunakan SQL untuk mengekstrak, membersihkan, dan mengolah data dari berbagai sumber
\item Memberikan rekomendasi berbasis data kepada tim manajemen terkait keputusan produk dan operasional
\end{itemize}
\textbf{Kata Kunci}: Data Visualization, SQL, Tableau, Cohort Analysis, Segmentasi Pelanggan, Excel}
\section{Bahasa}

\cvline{Indonesia}{Bahasa Ibu / Bilingual}
\cvline{Inggris}{Tingkat Profesional Penuh \hfill \textbf{Kata Kunci}: TOEFL iBT 95}
\section{Keahlian}

\cvline{Machine Learning}{Ahli \hfill \textbf{Kata Kunci}: Python, Scikit-learn, XGBoost, TensorFlow, Time Series, Model Interpretability}
\cvline{Data Analytics}{Ahli \hfill \textbf{Kata Kunci}: SQL, Pandas, NumPy, Statistics, A/B Testing, Feature Engineering}
\cvline{Data Engineering \& Cloud}{Menengah \hfill \textbf{Kata Kunci}: Apache Airflow, dbt, Google BigQuery, Cloud Composer, Docker}
\cvline{Data Visualization}{Mahir \hfill \textbf{Kata Kunci}: Tableau, Looker Studio, Matplotlib, Seaborn, Streamlit}
\section{Penghargaan}

\cventry{Okt 2021}
        {PT Finansial Cerdas}
        {Finalis Kompetisi Data Science Fintech}
        {}
        {}
        {Menjadi salah satu finalis dalam kompetisi prediksi risiko kredit dengan pendekatan ensemble learning.}
\section{Sertifikat}

\cventry{Mar 2023}
        {TensorFlow}
        {TensorFlow Developer Certificate}
        {\url{https://www.tensorflow.org/certificate}}
        {}
        {}

\cventry{Agu 2022}
        {Google Cloud}
        {Google Cloud Professional Data Engineer}
        {\url{https://cloud.google.com/certification/data-engineer}}
        {}
        {}
\section{Publikasi}

\cventry{Jul 2023}
        {Penerapan Gradient Boosting untuk Prediksi Churn Pelanggan pada Platform E-Commerce}
        {Jurnal Teknologi Informasi dan Ilmu Komputer}
        {\url{https://jtiik.ub.ac.id/index.php/jtiik}}
        {}
        {Publikasi penelitian tentang perbandingan performa beberapa algoritma machine learning dalam memprediksi churn pelanggan pada data transaksi e-commerce.}
\section{Proyek}

\cventry{Jan 2023–Jun 2023}
        {Mengembangkan sistem rekomendasi berbasis collaborative filtering dan content-based filtering untuk platform e-commerce.
}
        {Sistem Rekomendasi Produk}
        {\url{https://github.com/rakamahendra/recommendation-system}}
        {}
        {\begin{itemize}
\item Mengolah data interaksi pengguna dan metadata produk untuk membangun fitur rekomendasi
\item Mengevaluasi model dengan metrik precision@k dan recall@k
\item Mengimplementasikan model sebagai REST API menggunakan FastAPI
\end{itemize}
\textbf{Kata Kunci}: Recommendation System, Python, Collaborative Filtering, FastAPI}

\cventry{Feb 2022–Mei 2022}
        {Membangun dashboard interaktif untuk memantau performa penjualan dan inventaris.
}
        {Dashboard Analisis Penjualan}
        {\url{https://github.com/rakamahendra/sales-dashboard}}
        {}
        {\begin{itemize}
\item Menggunakan SQL untuk transformasi data dari database transaksional ke skema star
\item Menyusun visualisasi KPI utama seperti revenue, conversion rate, dan retensi
\item Menyajikan dashboard melalui Looker Studio untuk tim manajemen
\end{itemize}
\textbf{Kata Kunci}: SQL, Looker Studio, Data Visualization, Dashboard}
\section{Minat}

\cvline{Data dan Teknologi}{Machine Learning, Open Source, Data Journalism}
\cvline{Olahraga}{Lari, Sepak Bola, Badminton}
\cvline{Membaca}{Fiksi Sains, Biografi}
\section{Kegiatan Sukarela}

\cventry{Jan 2023–Sekarang}
        {Mentor Data Science}
        {Indonesia Data Science Society}
        {\url{https://www.idsc.io}}
        {}
        {\begin{itemize}
\item Membimbing peserta bootcamp dalam memahami konsep dasar machine learning dan statistika
\item Menyusun materi pelatihan dan studi kasus untuk peserta tingkat pemula
\item Memberikan umpan balik yang konstruktif pada proyek akhir peserta
\end{itemize}}

\end{document}